\chapter{Magnetic relaxation and QTM (menu 36)}
\label{ch:menu36}
\menulabel{36}

\section{Purpose}

A lanthanide single-molecule magnet is characterised by the barriers and
tunnelling processes its ab initio levels imply.  This menu reads what an
ORCA calculation already contains and organises it against the practical
guide's criteria: the per-doublet g-tensors and their mutual orientation
(the empirical effective-barrier estimate), the UBAR magnetic-moment matrix
elements (the qualitative relaxation path the engine itself prints), the
spin-orbit level structure, and -- when the Orca\_Magrelax section is
present -- the relaxation-rate table $\tau(T)$ with an Arrhenius/Orbach fit.
Nothing here is a dynamics simulation; every threshold carries its source.

\section{What you need}

One ORCA 6.1 output file carrying either

\begin{itemize}
  \item the \textbf{SINGLE\_ANISO} embedded section -- the \code{ANISO}
    sub-block inside \code{\%casscf} with \code{doaniso true} (manual
    section 5.32); give \code{MLTP} explicitly, otherwise the engine prints
    no g/D analysis at all (measured);
  \item the \textbf{Orca\_Magrelax} section -- \code{projectHSOC true} plus
    \code{projectedstates} in the CASSCF block, \code{DoSOC}/\code{DoMagrelax}
    in its \code{rel} sub-block, \code{\%qgprop} (numerical derivatives of
    the QDPT matrix along the normal modes) and \code{\%magrelax} (manual
    section 5.30).
\end{itemize}

A file with both gets both parts of the report.

\section{Walkthrough}

The example runs both parts on the CO$^+$ fixtures: first the
SINGLE\_ANISO output with \code{MLTP 2,2} (two Kramers doublets), then the
Orca\_Magrelax run of the same system -- a single vibrational mode at
2299.9~cm$^{-1}$ against a 14680.6~cm$^{-1}$ Zeeman gap, whose rate
table comes out all zeros (the measured boundary case).

\lstinputlisting[style=fbkcode, firstline=54, lastline=74,
  title={The menu-36 example (the SINGLE\_ANISO part: the per-KD table with
  the guard marks, the barrier estimate and the UBAR summary)}]
  {examples/expected/m36_relaxation.txt}

\section{What you get}

The report file \code{<output>.relax.fbk.md} carries:

\begin{itemize}
  \item per pseudospin group: energy, g principal values, the angle
    $\theta$ between the group's largest-g axis and the ground group's, and
    the guide's empirical average transverse g
    $g_T = (g_1 + g_2 + g_3 \sin\theta_3)/3$ with the product $g_T\theta$;
  \item the barrier estimate: the first excited group flagged by
    $\theta > 15^\circ$ (non-collinear) or $g_T\theta > 20$ (the guide's
    rule from 20 Dy(III) SMMs) -- or the explicit statement that no
    computed group qualifies, with the bound given by the highest computed
    level;
  \item the UBAR block: Zeeman eigenstates and the magnetic-moment matrix
    elements, averaged as the engine prints them,
    $(\lvert\mu_X\rvert + \lvert\mu_Y\rvert + \lvert\mu_Z\rvert)/3$;
  \item the magrelax part: the $\tau(T)$ table and the Orbach fit
    $\tau = \tau_0 \exp(U_{\mathrm{eff}}/k_BT)$ (U$_{\mathrm{eff}}$ in K
    and cm$^{-1}$, $\tau_0$ in seconds, $R^2$).
\end{itemize}

Beside the report, the magrelax part writes a plot-ready companion,
\code{<output>.relax.fbk.csv}: the $\tau(T)$ table as plain rows
(\code{temperature\_K,rate\_per\_s,tau\_s}), so an Arrhenius figure can be
made in any plotting tool from the same numbers the report prints
(Chapter~\ref{ch:formats}).

\section{Conventions, cross-checks and boundaries (measured)}

\begin{readbox}
\begin{itemize}
  \item \textbf{The axis-angle criteria need a defined axis.}  Below a
    relative g spread of 1\,\% the main-axis direction is decided by the
    print's digits: such groups are marked and excluded from the flags
    (measured on CO$^+$, $\Delta g/g \sim 2\times10^{-4}$).  The 15$^\circ$
    and $g_T\theta > 20$ thresholds are empirical (the guide's own
    working values); the flags bound a plausible barrier, not a measured
    one.
  \item \textbf{Two segments when DoSSC is on}: the output then contains
    two complete SINGLE\_ANISO runs (SOC-only and SOC+SSC; O$_2$: D =
    2.175287 vs 3.185964 -- the same numbers the QDPT ZFS blocks print).
    Compare spectra before reading across segments.
  \item \textbf{The UBAR table is qualitative by the engine's own
    statement} ("only a qualitative relaxation path"); its "even number
    of electrons" sentence is a fixed template warning, printed for
    odd-electron systems too (the CO$^+$ fixture is odd-electron).
  \item \textbf{An all-zero rate table is a structure fact.}  The
    one-phonon (Orbach) channel needs a phonon energy matching a Zeeman
    gap; a single-mode toy has none (CO$^+$: 2299.9 vs
    1.4--1.5$\times10^4$~cm$^{-1}$).  The report says so instead of
    fitting.
  \item \textbf{Definitions differ across exits.}  The SINGLE\_ANISO
    pseudospin g-values and the QDPT effective-Hamiltonian g-matrix
    (menu 35's reader) are different objects -- close numbers, not
    identical ones; state which exit a number comes from.
  \item \textbf{Boundaries}: the magnetic-dilution $\tau_{QT}$ model
    (dipolar-field statistics; the dilution design rules) is implemented
    in the tunnelling-prediction menu (menu~\menu{46}): its dilution
    variant (probabilistic neighbours, seeded repeats at the median)
    sits beside the equivalent-Zeeman and spin-dipole routes there;
    polynuclear exchange and POLY\_ANISO belong to item 5.6; the full
    spin-phonon treatment is what Orca\_Magrelax itself approximates
    (its Raman-I term is averaged to zero by the implementation's own
    statement, manual section 5.30).
\end{itemize}
\end{readbox}

\section{The ORCA side}

The SINGLE\_ANISO interface generates \code{<inp>.CASSCF.anisofile} and
\code{<inp>.anisoinp} and calls \code{otool\_single\_aniso} automatically; a
Kramers doublet prints no ZFS/D block, an integer pseudospin adds the $m=0$
column to the Zeeman-eigenstate table, and the UBAR matrix elements come in
intra-group ("opposite magnetization") and inter-group
("neighbouring multiplets") tables.  The Orca\_Magrelax module prints its
own editable input (\code{<inp>.magrelaxinp}), reads the QGPROP derivatives
(\code{socder.magrelax}) and the \code{.hess} file, and tabulates
$1/\tau$ and $\tau$ per temperature.  Both sections were measured on ORCA
6.1.1 (fixtures \code{single\_aniso/} and \code{magrelax/}; both sections'
format notes are recorded with those fixtures).
